\section{课程学习总结与反思}

本次课程报告从基础命令、Shell 脚本和综合项目三个层次展开。基础命令部分主要帮助我熟悉 Linux 的日常操作方式，包括目录切换、文件管理、内容查看、权限修改、系统信息查看和管道组合；Shell 编程部分则把这些命令组织成一个可以交互运行的答题系统，进一步练习了条件判断、循环、函数、文件读写和调试方法；TermSync 协同编辑器则把课程中的网络编程、线程同步、进程间通信、终端界面和文件持久化综合到一个较完整的应用中。

从学习过程来看，Linux 程序设计并不只是记忆命令格式，更重要的是理解系统资源如何被组织和访问。文件可以通过路径、权限和重定向进行管理；进程可以通过命令和脚本被启动、观察和调试；多个简单工具可以通过管道组合成复杂流程；网络服务则需要考虑连接、并发、状态一致性和异常处理。TermSync 项目的实现让我更直观地理解了这些知识之间的联系：一个看似简单的终端编辑器，背后同时涉及 socket 通信、线程安全、消息协议、自动保存、日志记录和测试验证。

本次报告也暴露出一些可以继续改进的地方。基础命令部分还可以增加更多关于错误输出、重定向和权限数字表示法的实验；Shell 答题系统还可以补充随机抽题、题库格式校验和错题回顾；TermSync 项目还可以继续完善身份认证、异常重连、批量操作和更丰富的集成测试。后续如果继续扩展，我会优先从稳定性和可维护性入手，让程序不仅能够完成演示，也能更接近真实 Linux 环境中的长期运行服务。

总体来说，这次课程报告把“命令使用”“脚本编写”和“系统程序设计”串联在一起，使我对 Linux 的理解从单个命令逐渐过渡到完整应用。通过在服务器上部署、运行、截图和测试，我也更加熟悉了远程 Linux 环境中的开发流程。
